\section{Algorithm Consolidation: Frozen Method Definition}
\label{sec:algorithm-consolidation}

The experimental program through Experiment~15A has reached a point at which further open-ended mechanism changes would make it harder, rather than easier, to determine what the actual scientific contribution is. The purpose of this second part of the report is therefore different from Part~I. Instead of introducing another candidate mechanism and testing whether it works, the goal is to freeze the simplest mechanism supported by the accumulated evidence, determine its exact relationship to existing communication-efficient federated methods, and close the systems gaps required for a complete federated-learning algorithm.

The consolidation phase is organized around three questions:
\begin{enumerate}
    \item \textbf{Distinctiveness:} Is the final stateful event encoder mathematically distinct from the closest existing sign, error-feedback, sparsification, and event-triggered methods?
    \item \textbf{Protocol completeness:} Does the communication advantage survive when both uplink and downlink communication are accounted for in an explicit server--client synchronization protocol?
    \item \textbf{FL compatibility:} Does the same event encoder remain effective when its input is a conventional multi-step local-training update as in FedAvg, rather than one stochastic-gradient completion at a time?
\end{enumerate}
A positive answer to all three would justify treating the present mechanism as the basis of a new communication-efficient FL algorithm. A negative answer to any one question would instead determine a more precise and narrower contribution statement.

\subsection{Frozen Algorithm V1}

The consolidation phase starts by freezing the communication mechanism. Let $u_i^r\in\R^d$ denote a generic local learning update produced by client $i$ at logical update index or round $r$. The central abstraction is intentionally broader than the gradient-specific formulation used in much of Part~I: $u_i^r$ may be a stochastic gradient contribution, a local model delta after several optimization steps, or another client-side update generated by a compatible local optimizer.

Each client maintains a persistent communication state $z_i^r\in\R^d$. The state update is
\begin{equation}
    z_i^{r+1}=\rho z_i^r+\alpha_i^r u_i^r,
    \label{eq:consolidated-state}
\end{equation}
where $\rho\in[0,1]$ is the evidence-memory factor and $\alpha_i^r$ is a protocol-dependent scaling coefficient. In the asynchronous FedSGD-like experiments of Part~I, $\alpha_i^r$ contains the client weight and compute-rate normalization.

A coordinate event is generated when
\begin{equation}
    |z_{ij}^{r+1}|\geq\Delta_j.
\end{equation}
The transmitted event is
\begin{equation}
    e=(j,s),
    \qquad
    s=\sign(z_{ij}^{r+1}),
\end{equation}
which contains a coordinate address and one sign bit. The corresponding communication-state coordinate is reset after transmission,
\begin{equation}
    z_{ij}^{r+1}\leftarrow0,
\end{equation}
and the server applies the sparse signed model jump
\begin{equation}
    w_j^+=w_j+q_r s.
    \label{eq:consolidated-server-jump}
\end{equation}
The event resolution $q_r>0$ is a design variable independent of the evidence-memory dynamics. Experiment~15A shows that this separation is essential: a persistent $\rho\approx1$ can remain useful while $q_r$ should become progressively finer during late-stage optimization.

The frozen V1 mechanism therefore consists of
\begin{equation}
\boxed{
\text{persistent temporal evidence}
+\text{coordinate first-passage events}
+\text{address/sign payload}
+\text{full reset}
+\text{annealed update resolution}.
}
\end{equation}

\subsection{What is deliberately excluded from V1}

The consolidation algorithm does \emph{not} include several mechanisms investigated earlier in the project. Homeostatic thresholds, strict coordinate competition, timing-derived jump rules, hard local or global descent certificates, layer-specific threshold normalization, filter-level events, clustering, and personalization are excluded because the experiments did not establish that they are required for the final empirical mechanism. This exclusion is scientifically important: V1 is not the union of all ideas tested in Part~I, but the smallest mechanism that survived the sequence of falsification and scaling experiments.

Adaptive event-resolution control is also not part of V1. Experiment~15A demonstrates that stronger manual annealing is necessary at longer horizons, but the adaptive replacement proposed as Experiment~15B has not yet been established. The consolidation phase therefore treats $q_r$ as a prescribed schedule and asks whether the resulting already-observed mechanism is sufficient to support a complete algorithmic claim.

\subsection{Communication interpretation}

The event encoder is best viewed as a \emph{stateful temporal compressor} rather than as a stand-alone optimization algorithm. The client-side state integrates learning information that was not yet considered sufficiently persistent or large to justify transmission. Unlike a memoryless sign compressor, communication is not mandatory at every local update. Unlike fixed-budget Top-$K$ sparsification, the number and identity of transmitted coordinates emerge from the first-passage process. Update magnitude is represented jointly by the fixed or scheduled event quantum and by event frequency over time.

This viewpoint motivates the remainder of Part~II. The novelty audit must determine whether this complete stateful compressor is already known under another terminology. The FedAvg study must determine whether the abstraction in \eqref{eq:consolidated-state} genuinely accepts local model deltas. The end-to-end protocol study must determine whether the sparse server evolution in \eqref{eq:consolidated-server-jump} can also support communication-efficient client synchronization.
